% !PRINTED 654
\weberpage{654}
\begin{center}
\S\ 184. \gesperrt{Basisform und Verzweigungsfunktional.}
\end{center}

Unter der Basisform des Körpers $\Omega$ wollen wir die Linearform
 verstehen:
\weq{1}{
\tau = \omega_1t_1 + \omega_2t_2 + \cdots + \omega_nt_n\,,
}
in der $t_1, t_2, \ldots, t_n$ die Funktionalvariablen sind, und deren Koeffizienten
 $\omega_1, \omega_2, \ldots, \omega_n$ eine Minimalbasis von $\Omega$ bilden. Diese
Form ist der größte gemeinschaftliche Teiler von $\omega_1, \omega_2, \ldots, \omega_n$
und folglich eine \gesperrt{Einheit.} Aus $\tau$ kann man alle ganzen
Funktionen in $\Omega$ ableiten, indem man für die Variablen $t_1,
t_2, \ldots, t_n$ ganze Funktionen in $Z$ setzt.

Setzen wir
\weq{2}{
F(t) = N(t - \tau) = t^n + A_1t^{n-1} + \cdots + A_{n-1}t + A_n\,,
}
so sind die $A_1, A_2, \ldots, A_n$ holomorphe ganze Funktionale in
$\overline{Z}$, und $\tau$ genügt der Gleichung $n$ten Grades:
\weq{3}{
F(\tau) = 0.
}

Ist $\pi$ ein Primfunktional und
\[
N_a(\pi) = z - c\,,
\]
so kann man nach \S\ 182, 6. die \gesperrt{Konstanten} $c_1, c_2, \ldots, c_n$ aus
den Kongruenzen
\weq{4}{
\omega_1 \equiv c_1, \ \omega_2 \equiv c_2, \ \ldots, \ \omega_n \equiv c_n \quad (\mathrm{mod}\ \pi)
}
bestimmen, und wenn man
\weq{5}{\tau_0 = c_1t_1 + c_2t_2 + \cdots + c_nt_n}
setzt, so ist $\tau_0$ ein \gesperrt{konstantes Funktional}, und $\tau - \tau_0$ ist
durch $\pi$ teilbar. Es ist nun zu beweisen:

3. \gesperrt{Das Primfunktional $\pi$ ist der größte gemeinschaftliche
 Teiler von $\tau - \tau_0$ und $z - c$:}
\weq{6}{
\pi = u(\tau - \tau_0) + v(z - c)\,,
}
\gesperrt{und wenn daher $(z - c)$ durch eine höhere als die
erste Potenz von $\pi$ teilbar ist, so ist $\tau - \tau_0$ nur durch
die erste Potenz von $\pi$ teilbar.}

Zunächst ergibt sich aus der Definition \S\ 181, daß
\weq{7}{
\tau - \tau_0 = t_1(\omega_1 - c_1) + t_2(\omega_2 - c_2) + \cdots + t_n(\omega_n - c_n)
}
der größte gemeinschaftliche Teiler von $(\omega_1 - c_1), (\omega_2 - c_2), \ldots,
(\omega_n - c_n)$ ist. Ist nun
\weq{8}{
\omega = x_1\omega_1 + x_2\omega_2 + \cdots x_n\omega_n \equiv 0 \quad (\mathrm{mod}\ \pi)\,,
% [sic] printed "x_2\omega_2 + \cdots x_n\omega_n" with no "+" after \cdots, the sum requires "+ \cdots + x_n\omega_n" as in (1), (5) and (7)
}
%CONT>>
